\documentclass[10pt,letterpaper]{article}
\usepackage[T1]{fontenc}
\usepackage[utf8]{inputenc}
\usepackage{lmodern}
\usepackage[margin=0.76in]{geometry}
\usepackage[table]{xcolor}
\usepackage{array,tabularx,longtable,booktabs}
\usepackage{enumitem}
\usepackage{amssymb,amsmath}
\usepackage{hyperref}
\usepackage{fancyhdr,lastpage}
\usepackage{microtype}
\usepackage{parskip}
\definecolor{cdblue}{HTML}{203957}
\definecolor{cdgray}{HTML}{586677}
\hypersetup{colorlinks=true,linkcolor=cdblue,urlcolor=cdblue,pdftitle={Laboratory 03: Unit, Integration and Contract Tests}}
\pagestyle{fancy}
\fancyhf{}
\lhead{\small\color{cdgray}Continuous Delivery Engineering}
\rhead{\small\color{cdgray}Implementation lab}
\cfoot{\thepage\ / \pageref{LastPage}}
\setlength{\headheight}{14pt}
\setlength{\emergencystretch}{2em}
\setlist[itemize]{leftmargin=*,itemsep=3pt,topsep=4pt}
\setlist[enumerate]{leftmargin=*,itemsep=3pt,topsep=4pt}
\begin{document}
\begin{center}
{\LARGE\bfseries\color{cdblue}Laboratory 03: Unit, Integration and Contract Tests}\\[7pt]
{\large Python unittest}\\[6pt]
{\small Reference edition: 2026-10-30 $\bullet$ Original educational and operational material}
\end{center}
\noindent\color{cdblue}\rule{\linewidth}{1.2pt}\color{black}
\vspace{0.4em}
\section{Learning objective and setup}
Goal: Build a tiered test strategy for the sample API. Prerequisites: Python unittest. Use the bundled examples only in a disposable repository or isolated AWS sandbox.


\section{Safety and isolation preflight}
\noindent $\square$ Confirm least-privileged identity and no long-lived secrets.\par\vspace{0.35em}
\noindent $\square$ Confirm sandbox scope, spending budget, and cleanup owner.\par\vspace{0.35em}
\noindent $\square$ Confirm no customer data and no destructive production action.\par\vspace{0.35em}

\section{Procedure}
\begin{enumerate}
\item Add two boundary tests and one injected failure case.
\item Create a consumer contract expectation and version it.
\item Classify checks as fast, acceptance or exploratory.
\end{enumerate}

\section{Expected results}
\begin{itemize}
\item Tests execute independently; no hidden network or wall-clock dependency.
\item Capture source revision, command outputs, job run identifiers and artifact checksums.
\item Record an intentional failure case and explain the failure condition.
\end{itemize}

\section{Evidence log}
\textbf{Repository/ref and execution timestamp}\par
\noindent\rule{\linewidth}{0.35pt}\par\vspace{0.7em}
\textbf{Commands, outputs and evidence links}\par
\noindent\rule{\linewidth}{0.35pt}\par\vspace{0.7em}
\noindent\rule{\linewidth}{0.35pt}\par\vspace{0.7em}
\noindent\rule{\linewidth}{0.35pt}\par\vspace{0.7em}
\textbf{Deviation from baseline}\par
\noindent\rule{\linewidth}{0.35pt}\par\vspace{0.7em}
\noindent\rule{\linewidth}{0.35pt}\par\vspace{0.7em}

\section{Troubleshooting decision tree}
\begin{itemize}
\item If authorization fails: inspect principal and trust conditions; do not broaden to wildcard access.
\item If artifact mismatch occurs: halt and rebuild in the trusted build stage; never relabel.
\item If a runtime health gate fails: pause progression, collect telemetry, invoke the applicable rollback.
\end{itemize}

\section{Completion assessment}
Pass only if the acceptance criterion is demonstrated, evidence is reproducible, and the negative case is observable. Rate correctness, security, repeatability, traceability and recovery on a 0--4 scale (pass at 16/20).



\end{document}
